\section{角关联、偏振和$\delta$反演}\label{sec:5-7}

在建立原子核能级纲图（第~5.6~节）之后，确定各能级的自旋宇称 $I^\pi$、鉴别电磁辐射的多极阶数（偶极 $L=1$ 或四极 $L=2$）以及提取电四极/磁偶极振幅混合比 $\delta(E2/M1)$，构成了在束 $\gamma$ 谱学向微观核结构模型（如转子模型、粒子-转子耦合与手征/摆动动力学）过渡的核心纽带。单纯的 $\gamma$ 跃迁能量与符合连接仅能提供能级拓扑连通性，无法独立锁定量子态的角动量与宇称。

在束重离子反应在自旋排列方向上具有天然的各向异性（见第~5.1~节）。出射 $\gamma$ 射线的空间角分布、双探测器方向角关联（DCO）以及康普顿线偏振度，携带着辐射场多极性与跃迁几率振幅干涉相位的直接信息。然而，由宏观角分布与偏振反演微观多极混合比 $\delta$ 是一类典型的非线性反问题。若对相位约定、退取向衰减、门选条件以及二次有理方程的多解特性缺乏严格界定，极易在谱学推断中产生两类典型偏差：其一，将特定几何下的 DCO 比值经验数值（如 1.0 或 0.5）当成普遍适用的多极性分界线，忽略门跃迁多极性与自旋序列的调制；其二，在角分布方程存在两个相容数学根时，依据理论模型的偏好人为舍弃其中一个物理自洽解分支，将理论假设冒充为独立实验判决。本节从统计取向张量与角分布几何展开出发，严格定义实多极混合比 $\delta$ 的相位相约，推导条件 DCO 强度比与康普顿散射偏振度方程，建立似然剖面下的双根反演与多解仲裁准则。

\subsection{角动量取向、辐射场几何与角分布理论}\label{subsec:5-7-1}

在束重离子完全熔合蒸发反应中，由于入射相对运动轨道角动量强烈偏向于垂直于束流轴的平面，初始复合核的磁子能级呈现高度取向的不均匀分布。在经历了相继的轻核子蒸发与统计非共振 $\gamma$ 级联后，待测母能级 $|i\rangle$（自旋为 $I_i$）依然保留着沿束流轴（定为量子化 $z$ 轴）的轴对称取向。

母态的自旋取向状态由统计取向张量（Orientation Tensor）$B_k(I_i)$ 描述（Hamilton 1975 经典专著规范~\cite{book-hamilton-electromagnetic-interaction}）：
\begin{equation}\label{eq:5-7-bk-def}
B_k(I_i) \equiv \sqrt{2I_i + 1} \sum_{M_i} (-1)^{I_i - M_i} \langle I_i M_i I_i (-M_i) \mid k 0 \rangle P(M_i),
\end{equation}
其中 $P(M_i)$ 为磁子能级 $M_i$ 的统计布居几率，且满足归一条件 $\sum_{M_i} P(M_i) = 1$。由 Clebsch-Gordan 系数正交归一性，张量零阶项 $B_0(I_i) \equiv 1$ 对任意归一布居分布 $P(M_i)$ 恒成立（并不限于各向同性分布）。沿束流轴的轴对称性仅保证方位角量子数 $q = 0$（无非对角取向张量）；仅当系统满足非定向排列（Alignment）对称性 $P(M_i) = P(-M_i)$ 时，奇阶取向张量才严格为零（$B_{2k+1} = 0$），使得级数仅保留偶数阶项 $k = 0, 2, 4$；若存在非对称极化（Polarization, $P(M_i) \ne P(-M_i)$），奇阶张量（如 $B_1$）亦将出现。若以初始复合核取向 $B_k(J_{\mathrm{CN}})$ 为基准，前级未观测级联退取向由去取向因子 $U_k = \prod_m U_k(m)$ 描述：
\begin{equation}\label{eq:5-7-deorientation}
B_k(I_i) = B_k(J_{\mathrm{CN}}) \prod_m U_k(m).
\end{equation}

从取向母态 $|i\rangle$ 退激至末态 $|f\rangle$（自旋 $I_f$）发射的单色 $\gamma$ 射线，其相对于束流轴夹角 $\theta$ 的相对角分布函数 $W(\theta)$ 展开为偶阶 Legendre 多项式序列：
\begin{equation}\label{eq:5-7-w-theta-general}
W(\theta) = \sum_{k=0, 2, 4} B_k(I_i) A_k(\delta) Q_k P_k(\cos\theta),
\end{equation}
注意式~\eqref{eq:5-7-w-theta-general} 中的 $B_k(I_i)$ 已包含前级退取向效应，绝不可在此展开中再次重复乘以相同的 $U_k$。
其中：
\begin{itemize}
\item $P_k(\cos\theta)$ 为标准第 $k$ 阶 Legendre 多项式（$P_0 = 1$，$P_2 = \frac{1}{2}(3\cos^2\theta - 1)$，$P_4 = \frac{1}{8}(35\cos^4\theta - 30\cos^2\theta + 3)$）；
\item $Q_k$ 为探测器对源张开有限立体角引起的几何衰减因子（Finite Solid-Angle Correction Factor），作为标量因子的简化表示仅对轴对称（如同轴圆柱或同轴圆对称探测器）严格成立（一般非轴对称几何须进行全立体角多极响应积分），满足 $Q_0 = 1$ 与 $|Q_k| \le 1$（并不限定恒满足 $0 < Q_k < 1$）；
\item $A_k(\delta)$ 为由跃迁多极类型与多极混合比决定的辐射角分布系数；
\item 截断阶数 $k \le 4$ 建立在出射辐射场多极性不超过四极（$L, L' \le 2$）的基础之上（由张量选择定则 $k \le \min(2I_i, 2L_{\max})$ 严格限定）。
\end{itemize}
在实验数据分析中，角分布通常拟合为常规 Legendre 展开形式：
\begin{equation}\label{eq:5-7-w-theta-fit}
W(\theta) = 1 + a_2 P_2(\cos\theta) + a_4 P_4(\cos\theta),
\end{equation}
其中表观展开系数满足 $a_k = \frac{B_k(I_i) U_k Q_k A_k(\delta)}{B_0(I_i) U_0 Q_0 A_0}$。

\subsection{多极混合比$\delta$的相位相约与二次有理展开}\label{subsec:5-7-2}

当跃迁在角动量三角定则与宇称守恒下允许低阶辐射多极 $L$ 与高阶辐射多极 $L' = L + 1$ 同时存在时（最常见的是磁偶极 $M1$ 与电四极 $E2$ 竞争跃迁），出射场呈现多极相干叠加。

必须对实多极混合比 $\delta$ 给出严格的无量纲辐射几率振幅比定义（与第~5.9~节式~\eqref{eq:5-9-delta-def} 保持一致，采用 Rose-Brink 规范体系~\cite{lange-kumar-hamilton-1982-multipole-admixtures, book-hamilton-electromagnetic-interaction}）：
\begin{equation}\label{eq:5-7-delta-def}
\delta(E2/M1) \equiv \frac{\gamma(E2)}{\gamma(M1)} = \frac{\sqrt{3}}{10} \frac{E_\gamma}{\hbar c} \frac{\langle I_f \| \mathcal{M}(E2) \| I_i \rangle}{\langle I_f \| \mathcal{M}_S(M1) \| I_i \rangle}.
\end{equation}
在此定义下：
\begin{itemize}
\item $\delta$ 是出射光子场的电四极振幅与磁偶极振幅之比，是严格的无量纲纯数；
\item 它不是裸约化矩阵元的简单相比，而是内含了跃迁动量因子 $q = E_\gamma / (\hbar c)$ 与相应的常数项，磁多极算符采用与电荷算符统一量纲的磁荷型算符约定 $\mathcal{M}_S(M1) \equiv \mathcal{M}_\mu(M1) / c$（含 $\mu_N/c$ 因子，与第~5.9~节一致）；
\item 初末态波函数的全同规范位相在比值中精确对消，$\delta$ 的符号由跃迁算符定义与物理矩阵元内禀结构共同锁定。
\end{itemize}

在文献转换中，必须注意不同历史相位相约的映射关系：
除 Rose-Brink（RB）约定外，核物理谱学还存在 Biedenharn-Rose（BR）约定与 Krane-Steffen（KS）约定。不同的相位约定可能导致实数 $\delta$ 的正负号发生翻转。\emph{若将文献中的 $\delta$ 符号转换至另一体系，必须同时将角分布干涉项系数的定义进行同步转换；绝不能在保留原公式几何系数形式的同时单独强行翻转 $\delta$ 的正负号}。

对于纯 $L$ 辐射与 $L'=L+1$ 辐射的混合，辐射角分布系数 $A_k(\delta)$ 可严格表征为混合比 $\delta$ 的二次有理函数~\cite{lange-kumar-hamilton-1982-multipole-admixtures}：
\begin{equation}\label{eq:5-7-ak-rational}
A_k(\delta) = \frac{F_k(L L I_f I_i) + 2\delta F_k(L L' I_f I_i) + \delta^2 F_k(L' L' I_f I_i)}{1 + \delta^2},
\end{equation}
其中 $F_k(L L' I_f I_i)$ 为纯几何性质的 Ferentz-Rosenzweig 系数（由 Racah 代数与 Wigner $6j$ 符号精确确定，并已在 Hamilton 1975 附录中系统列表）。对于最核心的 $M1+E2$ 跃迁（$L=1, L'=2$），式~\eqref{eq:5-7-ak-rational} 表达为：
\begin{equation}\label{eq:5-7-ak-m1e2}
A_k(\delta) = \frac{F_k(1 1 I_f I_i) + 2\delta F_k(1 2 I_f I_i) + \delta^2 F_k(2 2 I_f I_i)}{1 + \delta^2}.
\end{equation}
分子中的交叉项 $2\delta F_k(1 2)$ 直观展示了电四极与磁偶极发射振幅之间的量子力学相干干涉。

\subsection{定向态方向角关联（DCO）与门谱几何修正}\label{subsec:5-7-3}

虽然一维角分布 $W(\theta)$ 理论优美，但在重离子在束实验中，直接测量单谱角分布往往受到靶核连续本底与杂质峰的严重干扰。在多探测器阵列中，更常用的实验方法是在二维符合矩阵（第~5.6~节）上建立\emph{定向态方向角关联}（Directional Correlation from Oriented States, DCO）。

考虑由级联跃迁 $\gamma_1$ 与 $\gamma_2$ 构成的符合事件。实验将探测器分别放置在特定的两个极角 $\theta_1$ 与 $\theta_2$（相对于束流轴，通常选取 $90^\circ$ 与前向/后向角度如 $30^\circ$ 或 $135^\circ$）。由非对称符合矩阵（第~5.6~节式~\eqref{eq:5-6-ordered-pairs}）导出的实验 DCO 比值严格定义为条件门选强度之比：
\begin{equation}\label{eq:5-7-dco-def}
R_{\mathrm{DCO}} \equiv \frac{N_{\gamma 1, \gamma 2}(\theta_1, \theta_2)}{N_{\gamma 1, \gamma 2}(\theta_2, \theta_1)} \times \frac{\varepsilon_{\mathrm{p}}(\theta_2, E_{\gamma 1}) \varepsilon_{\mathrm{p}}(\theta_1, E_{\gamma 2})}{\varepsilon_{\mathrm{p}}(\theta_1, E_{\gamma 1}) \varepsilon_{\mathrm{p}}(\theta_2, E_{\gamma 2})},
\end{equation}
其中 $\varepsilon_{\mathrm{p}}(\theta, E)$ 为该角度通道的绝对全能峰探测效率。

必须澄清针对 DCO 比值的方法学常见误区：
\begin{enumerate}
\item DCO 比值绝非单谱角分布比值 $W(\theta_1)/W(\theta_2)$：
门跃迁 $\gamma_2$ 的观测在量子力学上对中间核能级的磁子能级布居进行了再次条件过滤，改变了后续发射态的取向状态。
\item 不存在通用的“DCO 等于 1.0 即为四极、等于 0.5 即为偶极”的绝对分界：
理论 $R_{\mathrm{DCO}}$ 的数值强烈依赖于门跃迁 $\gamma_2$ 的多极性质与自旋变化：
\begin{itemize}
\item 若门跃迁 $\gamma_2$ 为拉伸四极跃迁（$\Delta I = 2$ 纯 $E2$），在典型几何（如 $\theta_1 = 30^\circ, \theta_2 = 90^\circ$）下，待测跃迁若是拉伸四极（$\Delta I = 2$ 纯 $E2$），$R_{\mathrm{DCO}} \approx 1.0$；若待测跃迁为纯偶极（$\Delta I = 1$），$R_{\mathrm{DCO}} \approx 0.6$ 左右；
\item 若门跃迁 $\gamma_2$ 本身改为纯偶极跃迁（$\Delta I = 1$），基准参考值发生倒置——此时纯四极待测跃迁的 $R_{\mathrm{DCO}} \approx 1.7$ 左右，而纯偶极待测跃迁的 $R_{\mathrm{DCO}} \approx 1.0$。
\end{itemize}
\item 剩余自旋取向宽度的影响：
实际重离子实验中母态取向受高斯分布宽度参数 $\sigma/I$ 调制。未知的级联退取向与侧向馈入会压低偏离各向同性的幅度，因此反演时必须将取向宽度纳入参数空间同步拟合。
\end{enumerate}

\subsection{康普顿散射偏振度与电磁宇称仲裁}\label{subsec:5-7-4}

角分布与 DCO 比值主要敏感于多极阶数 $L$ 与自旋变化 $\Delta I$，但在鉴别出射光子的电磁宇称性质（例如区分属于宇称翻转的电偶极 $E1$（宇称选择定则 $\pi_i \pi_f = (-1)^L = -1$）还是宇称不变的磁偶极 $M1$（$\pi_i \pi_f = (-1)^{L+1} = +1$））时通常不具备区分能力。此时必须测量 $\gamma$ 射线的线偏振度（Linear Polarization）。

利用 Clover 探测器中 4 块相互正交晶体构成的微型康普顿偏振计（第~5.2~节式~\eqref{eq:5-2-klein-nishina-pol}），光子发生康普顿散射的截面强烈偏向于垂直于电矢量 $\mathbf{E}$ 的方位平面。实验测得的散射不对称度为：
\begin{equation}\label{eq:5-7-asym-exp}
A_{\mathrm{exp}} = \frac{N_\perp - a(E_\gamma) N_\parallel}{N_\perp + a(E_\gamma) N_\parallel},
\end{equation}
其中 $N_\parallel$ 与 $N_\perp$ 分别为光子在平行于核反应平面与垂直于核反应平面内散射进入相邻晶体的计数，$a(E_\gamma) = N_\perp^{\mathrm{unpol}} / N_\parallel^{\mathrm{unpol}}$ 为利用无偏振放射源标定的仪器几何不对称校正系数。

线偏振度 $P$ 表达为实验不对称度与偏振灵敏度因数 $Q(E_\gamma)$ 之比：
\begin{equation}\label{eq:5-7-pol-degree}
P = \frac{A_{\mathrm{exp}}}{Q(E_\gamma)},
\end{equation}
其中 Clover 探测器的偏振灵敏度 $Q(E_\gamma)$ 依赖于光子能量与晶体间距（通常在几十至数百 keV 能区通过标准校准线标定）。

在取向核跃迁理论中，相对偏振度 $P(\theta)$ 满足经典展开：
依据 Hamilton 偏振度理论展开（Hamilton 1975 专著第~9~章规范~\cite{book-hamilton-electromagnetic-interaction}），纯电性多极辐射与纯磁性多极辐射在相同自旋序列下的线偏振度呈现符号反转：
\begin{equation}\label{eq:5-7-pol-theor}
P(\theta) = \frac{\sum_k \left[ B_k A_k^{\mathrm{pol}}(\delta) Q_k^{\mathrm{pol}} P_k^{(2)}(\cos\theta) \right]}{W(\theta)},
\end{equation}
其中偏振张量系数 $A_k^{\mathrm{pol}}(\delta)$ 由不同螺旋度态下的跃迁振幅确定。在纯电四极（$E2$）或纯电偶极（$E1$）与纯磁偶极（$M1$）极限下偏振符号具有明确对应；但对于 $M1+E2$ 混合跃迁，偏振度由两多极分量的干涉二次式确定，\emph{偏振度 $P$ 为正绝不等于直接证明跃迁必为纯 $E2$ 辐射}。混合态的偏振符号与幅度强烈依赖于具体角度 $\theta$、自旋变化 $\Delta I$ 以及混合比 $\delta$ 的取值。

\subsection{似然剖面、多解双根与物理分支仲裁}\label{subsec:5-7-5}

由于辐射角分布系数 $A_k(\delta)$（式~\eqref{eq:5-7-ak-m1e2}）与实验 DCO 比值 $R_{\mathrm{DCO}}(\delta)$ 关于多极混合比 $\delta$ 均为二次有理函数：
\begin{equation}\label{eq:5-7-dco-rational}
R_{\mathrm{DCO}}(\delta) = \frac{n_0 + n_1 \delta + n_2 \delta^2}{d_0 + d_1 \delta + d_2 \delta^2},
\end{equation}
其中展开系数由探测角几何、初末态自旋序列及门跃迁多极性由双光子角关联理论确定。将实测 $R_{\mathrm{DCO}}$ 代入后，得到关于 $\delta$ 的二次代数方程：
\begin{equation}\label{eq:5-7-quadratic-roots}
(n_2 - R_{\mathrm{DCO}} d_2) \delta^2 + (n_1 - R_{\mathrm{DCO}} d_1) \delta + (n_0 - R_{\mathrm{DCO}} d_0) = 0.
\end{equation}
依系数判别式与极值情况，该方程在实数域可呈现 0、1 或 2 个有限实数根分支，或呈现渐近极点（$\delta \to \pm\infty$）。当存在两个实根时，两根并不必然表现为一小一大，其拟合优度亦未必相同；偏振测量或内转换系数能够提供独立的鉴别条件，但在缺乏附加实验数据时，不可无故断言某一根已被排除。

为向读者直观展示这一二次反演多解特征及偏振的裁决机制，图~\ref{fig:5-7-angular-branch} 给出了规范教学模型的解析曲线与双根相交示意图。

\begin{figure}[!htbp]
\centering
\IfFileExists{Img/angular-branch-delta-likelihood.pdf}{%
  \includegraphics[width=0.92\textwidth]{Img/angular-branch-delta-likelihood.pdf}%
}{%
  \fbox{\parbox{0.92\textwidth}{\centering\vspace{2.5cm}（图 5.3 DCO 二次反演双根与偏振度破简并示意图：\texttt{Img/angular-branch-delta-likelihood.pdf}；展示 $R_{\mathrm{DCO}}$ 实验带与有理二次曲线相交产生的两支相容根，以及线性偏振度 $P$ 对多解的独立鉴别）\vspace{2.5cm}}}%
}
\caption{多极混合比 $\delta$ 的二次有理反演与偏振简并解除机制（规范教学模型）。(a) 展示由严格正定响应比值构造的 $R_{\mathrm{DCO}}(\delta)$ 关于 $\arctan\delta$ 的二次有理特征曲线，水平阴影带代表测量误差范围 $R_{\mathrm{DCO}} = 0.80 \pm 0.05$。实验带与曲线相交给出两个完全相容的数学根分支 $\delta_1 = -0.35$（$\arctan\delta_1 = -19.4^\circ$）与 $\delta_2 = +1.59$（$\arctan\delta_2 = +57.8^\circ$）；(b) 给出由双螺旋度正定光子发射响应构造的线偏振度 $P(\delta)$ 理论曲线。在两个根位置偏振度呈现显著差异（$P(\delta_1) \approx -0.55$ 表明磁性分量主导，$P(\delta_2) \approx +0.67$ 表明电性分量主导），直观说明仅靠 DCO 比值无法独立定解，必须引入康普顿散射偏振计或内转换电子谱方可破除双根简并。注意此处为自洽规范教学模型，非特定未知用户硬件实测拟合）。}
\label{fig:5-7-angular-branch}
\end{figure}

图~\ref{fig:5-7-angular-branch} 确立了低能核结构谱学分析的一条关键证据边界：
\begin{enumerate}
\item 独立实验数据的分支保留原则：
在仅有角分布或 DCO 比值测量、缺乏具有足够区分精度的偏振度或内转换系数数据时，\emph{必须在实验结果中如实保留两个互不相容但均满足统计拟合的 $\delta$ 分支}！
\item 理论模型偏好不得冒充独立实验判决：
在核结构理论比对中（例如轴对称刚性转子模型或三轴粒子-转子耦合模型），模型往往预言特定带间跃迁具有极小的 $\delta$ 或特定的四极主导性。研究人员绝不能为了迎合某种理论偏好，在实验报告中人为剔除另一个与数据相容的候选根。
\item 连续似然剖面（Likelihood Profile）构建：
在统计推断中，可通过构建轮廓似然比检验曲线 $-2\Delta\ln \max_{\boldsymbol{\eta}} \mathcal{L}(\delta, \boldsymbol{\eta}) + \text{const}$，在满足渐近正则性条件下近似对应于 $\Delta\chi^2(\delta)$ 剖面；若原始矩阵包含非泊松结构或系统协方差，基于已发表观测矢量 $\mathbf{y}$ 及其协方差矩阵 $\mathbf{V}$ 的广义最小二乘模型依然是完全有效且符合统计学公理的推断框架。由于 $\delta$ 在实数轴上可延伸至 $\pm \infty$，通常在有界坐标 $\theta_\delta \equiv \arctan\delta \in [-\pi/2, \pi/2]$ 上绘制全局拟合剖面，展示可能呈现的双极小值结构及其非对称置信区间，并将其作为后续第~5.10~节全篇测量模型整合的客观输入。
\end{enumerate}

以奇 $A$ 核 $^{163}\mathrm{Lu}$ 中三轴超形变摆动带（TSD）的研究为例（A. Görgen 等人 2004 年经典工作~\cite{campaign-g07-gorgen-2004-163lu-quadrupole-moments}），在反演 TSD2 向 TSD1 带间跃迁的绝对 $B(E2)_{\mathrm{out}}$ 与 $B(M1)_{\mathrm{out}}$ 几率时，作者直接借用了先前独立薄靶角关联与偏振实验测得的平均值 $\delta = -3.10^{+0.36}_{-0.44}$，并假定其在整个高自旋带保持稳定。必须澄清：该文自身并未在此厚靶高自旋测量中重新独立测定高自旋各态的偏振度，而是基于共享的先前实验输入；这一实例表明微观多极强度反演往往依赖于先验假定与外部辅助测量。